As 55.946 leituras das 31 coletas selecionadas foram mantidas na análise, sem subamostragem para igualar as classes. A validação \textit{leave-one-subject-out} (LOSO) reservou um dos dez vasos para teste em cada rodada. Os pesos das classes foram calculados exclusivamente nas leituras dos nove vasos de treinamento de cada rodada, pela expressão $w_c=N_{\mathrm{treino}}/(2n_{c,\mathrm{treino}})$, em que $n_{c,\mathrm{treino}}$ é o número de leituras da classe $c$ no treino. Esses pesos foram aplicados ao ajuste do \textit{Extra Trees}, \textit{Random Forest} e MLP. No k-NN, os votos dos cinco vizinhos foram ponderados pela distância e pelo peso de sua classe. Na TCN, a função de perda utilizou pesos calculados a partir das janelas dos vasos de treino. A acurácia balanceada e o F1 macro foram calculados nas previsões dos vasos reservados, sem reponderar o conjunto de teste. A unidade principal de avaliação foi a leitura; as previsões por coleta e por vaso foram obtidas pela média das probabilidades.
